\documentclass[11pt]{article}
\usepackage[margin=1.1in]{geometry}
\usepackage{amsmath,amssymb,amsthm}
\usepackage[colorlinks=true,linkcolor=blue,citecolor=blue]{hyperref}

\newtheorem{theorem}{Theorem}[section]
\newtheorem{lemma}[theorem]{Lemma}
\newtheorem{proposition}[theorem]{Proposition}
\newtheorem{corollary}[theorem]{Corollary}
\newtheorem{definition}[theorem]{Definition}
\newtheorem{remark}[theorem]{Remark}
\newtheorem{problem}[theorem]{Problem}

\newcommand{\Q}{Q_\infty}
\newcommand{\F}{\mathcal{F}}
\newcommand{\dist}{\operatorname{dist}}
\newcommand{\diam}{\operatorname{diam}}
\newcommand{\EX}{\mathsf{Ex}}

\title{Margin-stable schemes and the Baire category of full complexity\\
in non-autonomous dynamics on the Hilbert cube}
\author{Yitzchak Shmalo\\
\small Einstein Institute of Mathematics, The Hebrew University of Jerusalem, Jerusalem, Israel\\
\small }
\date{June 2026}

\begin{document}
\maketitle

\begin{abstract}
In the full non-autonomous mapping space $\F(\Q)=C(\Q,\Q)^{\mathbb N}$ of the Hilbert cube,
equipped with the uniform-in-time metric, the recent construction of
\cite{Shmalo2026} produces a \emph{dense} set of systems admitting exact windowed
Hilbert horseshoes, with all-tail symbolic semifactors, windowed distributional chaos of
type~1, infinite topological entropy, and infinite metric mean dimension; the same paper
observes that the exact class is not stable and asks (Section~14 of \cite{Shmalo2026})
for a residual theorem obtained by a ``Baire-fusion mechanism.''
This paper resolves the stability side of that question.  First, we introduce
\emph{strict depth-$N$ brick schemes}, margin-protected finite-coordinate covering
configurations, and prove that for every depth $N$ and every sequence of dense open
windows the corresponding class $\mathcal U_N$ is open and dense, with explicit
stability radius $\theta\,2^{-N-1}$ for a scheme of margin~$\theta$.  Consequently the
class $\mathcal S=\bigcap_N \mathcal U_N$ of systems carrying windowed strict schemes of
\emph{every} finite depth is residual.  Second, on $\mathcal S$ every branch brick covers
every depth-$N$ coordinate pattern of the next bricks \emph{exactly}, by a
Poincar\'e--Miranda argument; this is the maximal one-step covering compatible with
stability.  Third, a genuinely dynamical statement holds residually: for a residual
set of systems, every tail exhibits separated orbit growth of every order $q$ at
arbitrarily fine scales (``entropy bursts''), the residual shadow of the dense
infinite-metric-mean-dimension theorem.  Fourth, we prove an obstruction theorem: the
exact-horseshoe class of
\cite{Shmalo2026} has \emph{empty interior}, and indeed every property that forces some
image to have full range in arbitrarily deep coordinates has empty interior.  The
obstruction shows that the depth stratification is not an artifact: no margin-stable
strengthening of the exact theorem can be open, so the residual hierarchy
$\mathcal S$ is the natural generic object.  We close with a precise formulation of
what remains open.
\end{abstract}

\section{Introduction}\label{sec:intro}

A non-autonomous discrete dynamical system on a compact metric space $(X,\rho)$ is a
sequence $f=(f_1,f_2,\dots)$ of continuous self-maps of $X$, acting through the
compositions $f_n\circ\cdots\circ f_1$.  The full mapping space
\[
\F(X)=C(X,X)^{\mathbb N},
\qquad
D_X(f,g)=\sup_{n\ge 1}\,\sup_{x\in X}\rho\bigl(f_n(x),g_n(x)\bigr),
\]
is a complete, hence Baire, metric space (Lemma~\ref{lem:baire} below).  For interval
systems and for systems on finite-dimensional cubes and compact manifolds, the paper
\cite{Shmalo2026} constructs $D$-open dense sets of systems displaying \emph{all-tail
full complexity}: after deleting any finite initial segment, the remaining system
admits a compact symbolic semifactor onto a full variable-alphabet shift, a
continuum-sized scrambled set for distributional chaos of type 1 visiting prescribed
dense open windows at prescribed absolute times, infinite topological entropy, and the
sharp small-scale entropy spectrum $H_{g,r}(\varepsilon)=d\log(1/\varepsilon)+O(1)$.
The proofs rest on \emph{strict} cubical covering relations: face inequalities with a
positive margin $\theta$, which survive every perturbation of size less than $\theta$
and are converted into covering statements by the Poincar\'e--Miranda theorem
\cite{Miranda,Kulpa}.

On the Hilbert cube
\[
\Q=[0,1]^{\mathbb N},
\qquad
\rho(x,y)=\sum_{j=1}^\infty 2^{-j}\,|x_j-y_j|,
\]
the situation changes.  The result of \cite{Shmalo2026} on $\Q$ is a \emph{dense}
theorem: a dense set of systems admits exact windowed Hilbert horseshoes
(Definition~\ref{def:exact} below), built from finite-support boxes mapped
\emph{exactly} onto the next box.  Exactness in arbitrarily high coordinates is not
stable, because coordinate $j$ carries metric weight $2^{-j}$: a perturbation of
uniform size $\delta$ can move the $j$-th image coordinate by $2^{j}\delta$.  Section 14
of \cite{Shmalo2026} records the resulting question, which is the starting point of
the present paper:

\begin{quotation}
\noindent``The Hilbert-cube conclusion is dense but is not asserted to be open: exact
covering in arbitrarily high coordinates is not stable under the Hilbert metric.  A
residual Hilbert-cube theorem with coherent all-tail symbolic codings would require an
additional Baire-fusion mechanism beyond the exact finite-support implantations proved
here.''
\end{quotation}

This paper carries out the stability analysis that any such fusion mechanism must
confront, and proves four results: a residual depth hierarchy of margin-stable
schemes (Theorem~\ref{thm:A}), exact one-step covering at every depth on that
residual set (Theorem~\ref{thm:B}), a residual entropy-burst theorem in every tail
(Theorem~\ref{thm:burst}; see Section~\ref{sec:burst}), and an empty-interior
obstruction (Theorem~\ref{thm:C}).

\subsection*{Main results}

Fix a sequence $W_\bullet=(W_n)_{n\ge 1}$ of dense open subsets of $\Q$ (\emph{windows})
and geometric data $(M_n,q_n,\gamma_n)_{n\ge1}$ as in Section~\ref{sec:schemes}
(finitely many branches $M_n\ge 2$, separation dimensions $q_n\ge 1$, diameter bounds
$\gamma_n>0$; the block schedule of \cite[Section~12]{Shmalo2026} is the motivating
instance).  For each integer $N\ge 1$ we define in Section~\ref{sec:schemes} the class
\[
\mathcal U_N(W_\bullet)\subseteq\F(\Q)
\]
of systems admitting a \emph{strict depth-$N$ windowed brick scheme}: a time-indexed
family of pairwise separated finite-support \emph{bricks} inside the windows, such that
every brick crosses, in every coordinate $\ell\le N$ and with a uniform margin
$\theta>0$, the coordinate range of all bricks at the next time.

\begin{theorem}[Finite-depth genericity; Theorem~\ref{thm:main} below]\label{thm:A}
For every depth $N\ge1$ and every window sequence $W_\bullet$, the class
$\mathcal U_N(W_\bullet)$ is open and dense in $(\F(\Q),D)$.  A scheme with margin
$\theta$ persists, with margin $\theta/2$, on the entire ball of radius
$\theta\,2^{-N-1}$ around the realizing system.  Consequently
\[
\mathcal S(W_\bullet)\;=\;\bigcap_{N\ge1}\mathcal U_N(W_\bullet)
\]
is residual in $\F(\Q)$, as is the simultaneous class
$\bigcap_{a\ge1}\mathcal S(W^a_\bullet)$ for any countable family of window sequences.
\end{theorem}

\begin{theorem}[Generic exact one-step covering; Theorem~\ref{thm:covering} below]\label{thm:B}
Let $g\in\mathcal S(W_\bullet)$.  Then for every $N\ge 1$ there is a depth-$N$ scheme
for $g$ such that for every time $n$, every brick $B_{n,i}$ of the scheme, and every
target pattern $t=(t_\ell)_{\ell\le N}$ lying coordinatewise in the interval spanned
by the bricks at time $n+1$, there exists $x\in B_{n,i}$ with
\[
\pi_\ell\,g_n(x)=t_\ell\qquad(1\le\ell\le N).
\]
In particular every brick maps exactly onto every depth-$N$ coordinate pattern of every
brick at the next time, inside the prescribed windows.
\end{theorem}

\begin{theorem}[Obstruction: the exact class has empty interior; Theorem~\ref{thm:empty} below]\label{thm:C}
Let $\EX\subseteq\F(\Q)$ denote the set of systems admitting an exact windowed Hilbert
horseshoe in the sense of \cite[Definition~12.1]{Shmalo2026}
(Definition~\ref{def:exact} below).  Then $\EX$ is dense and has empty interior.
More generally, every class $\mathcal P\subseteq\F(\Q)$ with the property that each
$g\in\mathcal P$ has, for some $n$ and some compact $A\subseteq\Q$, an image $g_n(A)$
whose coordinate projections are onto $[0,1]$ in infinitely many coordinates, has empty
interior.
\end{theorem}

Theorem~\ref{thm:C} shows that the ``dense, not open'' phenomenon of
\cite{Shmalo2026} is intrinsic and not an artifact of the construction: \emph{no}
sharpening of exact infinite-depth covering can be open, since arbitrarily small
uniform perturbations (an explicit coordinate squeeze) destroy full range in deep
coordinates everywhere in the space.  Theorems~\ref{thm:A} and~\ref{thm:B} identify
the optimal margin-stable substitute: full strict covering of every finite depth,
simultaneously for all depths, on a residual set.  The stability radius
$\theta2^{-N-1}$ quantifies the cost of depth and explains the stratification: the
margins of conditions at depth $N$ have stability comparable to $2^{-N}$, so no single
open set can certify unbounded depth, while the intersection over $N$ retains
residuality.

\subsection*{Discussion: what the fusion mechanism can and cannot give}

It is natural to ask whether membership in $\mathcal S(W_\bullet)$ already implies the
\emph{dynamical} conclusions of the dense theorem of \cite{Shmalo2026} (semifactors,
DC1, infinite entropy and metric mean dimension in every tail).  We do not claim this,
and we believe the question is genuinely delicate, for a reason worth recording.  The
infinite-horizon conclusions all require orbits threading prescribed bricks at
\emph{all} times.  The backward-induction production of such orbits in
\cite{Shmalo2026} consumes exact surjectivity onto the next box in \emph{all}
coordinates; a depth-$N$ scheme controls only $N$ image coordinates per step, and the
uncontrolled deep coordinates of the image re-enter as uncontrolled \emph{inputs} at
the next step.  Quantitatively, a condition involving image coordinate $\ell$ has
stability $O(2^{-\ell})$ under $D$, while restoring such a condition after a
perturbation of size $r$ costs $\Omega(\min(1,2^{\ell}r))$ per coordinate in amplitude;
summing over the critical band of coordinates, maintenance of unbounded-depth margins
along a Baire fusion fails against an adversary that shrinks neighbourhoods faster
than geometrically.  This is exactly the obstruction isolated in
Theorem~\ref{thm:C} in its sharpest open form.  We therefore state the remaining
problem precisely in Section~\ref{sec:open}, in a form that separates what is now
known (the residual depth hierarchy, with exact one-step covering at every depth, plus
the empty-interior obstruction) from what is open (residuality of the infinite-horizon
conclusions themselves).

\subsection*{Relation to the finite-dimensional theory}

On $Q_d=[0,1]^d$ with the sup metric, all $d$ coordinates have unit weight, margins
have uniform stability, and \cite{Shmalo2026} obtains a single open dense class with
the full dynamical conclusions.  The present results show in what precise sense
$d=\infty$ differs: openness survives at every finite depth with an exponentially
small but positive radius, and fails at infinite depth for structural reasons.  The
residual set $\mathcal S(W_\bullet)$ is the direct analogue of the finite-dimensional
open dense class under this constraint.

\subsection*{Organization}

Section~\ref{sec:prelim} fixes notation and records the Baire property.
Section~\ref{sec:schemes} defines strict depth-$N$ brick schemes and proves openness
and the covering theorem.  Section~\ref{sec:density} proves the selection and
implantation lemmas and density.  Section~\ref{sec:residual} assembles
Theorem~\ref{thm:A}.  Section~\ref{sec:burst} proves the residual entropy-burst
theorem.  Section~\ref{sec:obstruction} proves the obstruction theorem.
Section~\ref{sec:open} states the remaining open problem.

\section{Preliminaries}\label{sec:prelim}

Throughout, $\Q=[0,1]^{\mathbb N}$ carries the metric
$\rho(x,y)=\sum_{j\ge1}2^{-j}|x_j-y_j|$, which induces the product topology and makes
$\Q$ compact.  We write $\pi_\ell:\Q\to[0,1]$ for the $\ell$-th coordinate projection.
The mapping space $C(\Q,\Q)$ carries the uniform metric
$d_\infty(u,v)=\sup_{x}\rho(u(x),v(x))$, and
\[
\F(\Q)=C(\Q,\Q)^{\mathbb N},\qquad
D(f,g)=\sup_{n\ge1}d_\infty(f_n,g_n).
\]

\begin{lemma}[{\cite[Lemma 2.1]{Shmalo2026}}]\label{lem:baire}
$(\F(\Q),D)$ is a complete metric space, and hence a Baire space.
\end{lemma}

\begin{proof}
$C(\Q,\Q)$ is complete because $\Q$ is compact and uniform limits of continuous maps
are continuous.  If $(f^{(k)})_k$ is $D$-Cauchy, each coordinate sequence
$(f^{(k)}_n)_k$ is $d_\infty$-Cauchy with a modulus independent of $n$; the coordinate
limits $f_n$ assemble to $f\in\F(\Q)$ with $D(f^{(k)},f)\to0$.  Completeness implies
the Baire property.
\end{proof}

We use the elementary weight estimate repeatedly: for $u,v\in C(\Q,\Q)$ and any
coordinate $\ell$,
\begin{equation}\label{eq:weight}
\sup_{x\in\Q}\,\bigl|\pi_\ell u(x)-\pi_\ell v(x)\bigr|
\;\le\;2^{\ell}\,d_\infty(u,v),
\qquad
d_\infty(u,v)\;\le\;\sum_{\ell\ge1}2^{-\ell}\sup_x|\pi_\ell u(x)-\pi_\ell v(x)|.
\end{equation}

A \emph{finite-support box} is a product $B=\prod_{j\ge1}I_j$ with each
$I_j\subseteq[0,1]$ a nondegenerate closed interval and $I_j=[0,1]$ for all $j>R$; the
least such $R$ is the \emph{support level}.  Finite-support boxes are compact with
nonempty interior, and they form a local base: every nonempty open $U\subseteq\Q$
contains finite-support boxes of arbitrarily small diameter
\cite[Lemma~10.2]{Shmalo2026}.  For $K\ge1$ and $\alpha\in(0,1)$, an
\emph{$(\alpha,K)$-box} is a finite-support box of the form
$J=\prod_{j=1}^{K}[a_j,a_j+\alpha]\times[0,1]^{\{K+1,K+2,\dots\}}$; then
$\diam_\rho(J)\le\alpha+2^{-K}$.

Finally we recall the Poincar\'e--Miranda theorem \cite{Miranda,Kulpa}.

\begin{theorem}[Poincar\'e--Miranda]\label{thm:miranda}
Let $P=\prod_{\ell=1}^{N}[r_\ell,s_\ell]\subset\mathbb R^N$ and let
$\Phi=(\Phi_1,\dots,\Phi_N):P\to\mathbb R^N$ be continuous.  If for every $\ell$,
$\Phi_\ell\le0$ on $\{z\in P: z_\ell=r_\ell\}$ and $\Phi_\ell\ge0$ on
$\{z\in P: z_\ell=s_\ell\}$, then $\Phi$ has a zero in $P$.
\end{theorem}

\section{Strict depth-\texorpdfstring{$N$}{N} brick schemes}\label{sec:schemes}

Fix once and for all \emph{geometric data}: integers $M_n\ge2$ (branch counts),
$q_n\ge1$ (separation dimensions), and reals $\gamma_n>0$ (diameter bounds), $n\ge1$.
For the dynamical applications of \cite{Shmalo2026} one takes the universal block
schedule of \cite[Section~12]{Shmalo2026}, where every label $C_a$, $D_s$, $E_{q,m}$
occurs on infinitely many blocks and $\gamma_n\to0$; none of the results below depends
on that specific choice.  Fix also a sequence $W_\bullet=(W_n)$ of dense open subsets
of $\Q$.

\begin{definition}[Bricks]\label{def:brick}
A \emph{brick} with support level $R$ is a finite-support box
$B=\prod_{j\le R}I_j\times[0,1]^{\{R+1,\dots\}}$ with all restricted intervals $I_j$
nondegenerate.  For bricks $B,B'$ and $q\ge1$ we write
\[
\dist_q(B,B')=\inf_{x\in B,\,y\in B'}\ \max_{1\le\ell\le q}|x_\ell-y_\ell|.
\]
\end{definition}

\begin{definition}[Strict depth-$N$ windowed brick scheme]\label{def:scheme}
Let $N\ge1$.  A system $g\in\F(\Q)$ \emph{admits a strict depth-$N$
$W_\bullet$-windowed brick scheme} if there exist $\alpha\in(0,1/4)$, a margin
$\theta\in(0,1/16]$, support levels $R_n\ge\max(N,q_n)$, and bricks
\[
B_{n,1},\dots,B_{n,M_n}\subseteq W_n,
\qquad
B_{n,i}=\prod_{j\le R_n}I^{n,i}_j\times[0,1]^{\{R_n+1,\dots\}},
\]
such that, writing
\[
m^-_{n,\ell}=\min_{1\le i\le M_n}\min I^{n,i}_\ell,
\qquad
m^+_{n,\ell}=\max_{1\le i\le M_n}\max I^{n,i}_\ell
\qquad(\ell\le N),
\]
the following hold for all $n\ge1$:
\begin{enumerate}
\item[(S0)] (\emph{anchoring}) $I^{n,i}_\ell\subseteq[2\theta,\,1-2\theta]$ for all
$i$ and all $\ell\le N$;
\item[(S1)] $\diam_\rho(B_{n,i})<\gamma_n$ for all $i$;
\item[(S2)] $\dist_{q_n}(B_{n,i},B_{n,j})>c_{\alpha,q_n}M_n^{-1/q_n}$ for $i\ne j$,
where $c_{\alpha,q}>0$ are the selection constants of
Lemma~\ref{lem:selection};
\item[(S3)] (\emph{strict depth-$N$ crossing}) for every $i$ and every
$\ell\le N$,
\[
\pi_\ell\,g_n(x)\;\le\;m^-_{n+1,\ell}-\theta
\quad\text{for all }x\in B_{n,i}\text{ with }x_\ell=\min I^{n,i}_\ell,
\]
\[
\pi_\ell\,g_n(x)\;\ge\;m^+_{n+1,\ell}+\theta
\quad\text{for all }x\in B_{n,i}\text{ with }x_\ell=\max I^{n,i}_\ell .
\]
\end{enumerate}
We write $\mathcal U_N(W_\bullet)$ for the set of systems admitting such a scheme.
By (S0), $m^-_{n+1,\ell}-\theta\ge\theta>0$ and $m^+_{n+1,\ell}+\theta\le1-\theta<1$,
so the inequalities in (S3) are never vacuous for range reasons.
\end{definition}

\begin{lemma}[Openness with quantified radius]\label{lem:open}
Let $g\in\mathcal U_N(W_\bullet)$, realized by a scheme with margin $\theta$.  If
$h\in\F(\Q)$ satisfies $D(h,g)<\theta\,2^{-N-1}$, then $h$ realizes the \emph{same}
bricks as a strict depth-$N$ scheme with margin $\theta/2$.  In particular
$\mathcal U_N(W_\bullet)$ is open in $(\F(\Q),D)$.
\end{lemma}

\begin{proof}
Conditions (S1) and (S2) and the window membership concern only the bricks.  For (S3),
fix $n,i,\ell\le N$ and $x\in B_{n,i}$ with $x_\ell=\min I^{n,i}_\ell$.  By
\eqref{eq:weight},
\[
\pi_\ell h_n(x)\;\le\;\pi_\ell g_n(x)+2^{\ell}D(h,g)
\;<\;\bigl(m^-_{n+1,\ell}-\theta\bigr)+2^{N}\,\theta\,2^{-N-1}
\;=\;m^-_{n+1,\ell}-\tfrac{\theta}{2},
\]
and symmetrically on the faces $x_\ell=\max I^{n,i}_\ell$.
\end{proof}

\begin{theorem}[Exact one-step depth-$N$ covering]\label{thm:covering}
Let $g$ admit a strict depth-$N$ scheme with bricks $B_{n,i}$ and margin $\theta$.
Then for every $n\ge1$, every $1\le i\le M_n$, and every target pattern
$t=(t_\ell)_{\ell\le N}$ with
$t_\ell\in[m^-_{n+1,\ell},\,m^+_{n+1,\ell}]$ for all $\ell\le N$,
there exists $x\in B_{n,i}$ with
\[
\pi_\ell\,g_n(x)=t_\ell\qquad(1\le\ell\le N).
\]
In particular, for every $1\le i'\le M_{n+1}$ there is $x\in B_{n,i}$ whose image
$g_n(x)$ agrees in its first $N$ coordinates with a prescribed point of $B_{n+1,i'}$;
all such $x$ lie in $B_{n,i}\subseteq W_n$.
\end{theorem}

\begin{proof}
Write $I_\ell=I^{n,i}_\ell$.  Define the section
$\sigma:\prod_{\ell\le N}I_\ell\to B_{n,i}$ by
$\sigma(z)_\ell=z_\ell$ for $\ell\le N$, $\sigma(z)_j=\operatorname{mid}I^{n,i}_j$ for
$N<j\le R_n$, and $\sigma(z)_j=1/2$ for $j>R_n$; $\sigma$ is continuous and takes
values in $B_{n,i}$.  Let
\[
\Phi_\ell(z)=\pi_\ell\,g_n(\sigma(z))-t_\ell\qquad(\ell\le N).
\]
If $z_\ell=\min I_\ell$, then $\sigma(z)\in B_{n,i}$ has $\ell$-th coordinate
$\min I_\ell$, so by (S3),
$\Phi_\ell(z)\le m^-_{n+1,\ell}-\theta-t_\ell\le-\theta<0$; if
$z_\ell=\max I_\ell$, then $\Phi_\ell(z)\ge\theta>0$.  By
Theorem~\ref{thm:miranda} there is $z$ with $\Phi(z)=0$, and $x=\sigma(z)$ is the
required point.  For the final assertion take
$t_\ell\in I^{n+1,i'}_\ell\subseteq[m^-_{n+1,\ell},m^+_{n+1,\ell}]$ for $\ell\le N$.
\end{proof}

\begin{remark}\label{rem:depth-cost}
The radius $\theta2^{-N-1}$ in Lemma~\ref{lem:open} is the correct order of the
stability of depth-$N$ information: a perturbation supported in the $N$-th image
coordinate of size $2\theta$ in amplitude has $D$-size $2^{-N+1}\theta$ and destroys
the margin-$\theta$ inequalities (S3) at $\ell=N$ for the given scheme.  Stability of
margin conditions thus decays exponentially in depth, which is the quantitative form
of the instability observed in \cite[Section~14]{Shmalo2026} and the reason the
classes $\mathcal U_N$ are stratified by depth rather than united in a single open
class; Theorem~\ref{thm:empty} below makes the impossibility of the united class
precise.
\end{remark}

\section{Selection, implantation, and density}\label{sec:density}

We first record the measure-theoretic separation lemma and the windowed selection
lemma of \cite{Shmalo2026}, in the form needed here.

\begin{lemma}[{Separated points; \cite[Lemma 5.1]{Shmalo2026}}]\label{lem:separated}
Let $A\subseteq\mathbb R^{q}$ be bounded and measurable with $\lambda_q(A)>0$.  For
every $m\ge2$ the set $A$ contains $m$ points whose pairwise sup distances exceed
$\tfrac14\,\lambda_q(A)^{1/q}m^{-1/q}$.
\end{lemma}

\begin{proof}
Put $r=\tfrac1{16}(\lambda_q(A)/m)^{1/q}$ and let $P\subseteq A$ be a maximal
$4r$-separated subset.  If $\operatorname{card}(P)\le m-1$, maximality lets $A$ be
covered by at most $m-1$ sup-balls of radius $4r$, whence
$\lambda_q(A)\le(m-1)(8r)^q<2^{-q}\lambda_q(A)$, a contradiction.
\end{proof}

\begin{lemma}[Windowed selection]\label{lem:selection}
Fix $\alpha\in(0,1/4)$ and $q\ge1$.  There is a constant $c_{\alpha,q}>0$ with the
following property.  Let $J=\prod_{j\le K}[u_j,u_j+\alpha]\times[0,1]^{\{K+1,\dots\}}$
be an $(\alpha,K)$-box with $K\ge q$, let $W\subseteq\Q$ be dense open, let
$f\in C(\Q,\Q)$, and let $m\ge2$.  Then there are points
$p_1,\dots,p_m\in\operatorname{int}(J)\cap W$ such that
\[
\max_{1\le\ell\le q}\bigl|(p_i)_\ell-(p_j)_\ell\bigr|>4c_{\alpha,q}\,m^{-1/q}
\quad(i\ne j),
\qquad
\diam_\rho\{f(p_1),\dots,f(p_m)\}<\alpha/16 .
\]
\end{lemma}

\begin{proof}
Define a slice $\iota:P\to\operatorname{int}(J)$,
$P=\prod_{\ell\le q}[u_\ell+\alpha/4,\,u_\ell+3\alpha/4]$, by
$(\iota z)_\ell=z_\ell$ for $\ell\le q$, $(\iota z)_j=u_j+\alpha/2$ for $q<j\le K$,
and $(\iota z)_j=1/2$ for $j>K$.  By compactness cover $\Q$ by Borel sets
$R_1,\dots,R_P$ of $\rho$-diameter less than $\alpha/64$, where $P=P_\alpha$ depends
only on $\alpha$.  The sets $A_s=\{z\in P: f(\iota z)\in R_s\}$ cover $P$, so some
$A_s$ has $\lambda_q(A_s)\ge(\alpha/2)^qP_\alpha^{-1}$.  Lemma~\ref{lem:separated}
gives $z_1,\dots,z_m\in A_s$ pairwise more than $8c_{\alpha,q}m^{-1/q}$ apart in sup
norm, where $c_{\alpha,q}=\tfrac1{32}\bigl((\alpha/2)^qP_\alpha^{-1}\bigr)^{1/q}$.
All $f(\iota z_i)$ lie in one $R_s$.  By continuity of $f$, choose pairwise disjoint
relatively open $U_i\ni\iota z_i$ inside $\operatorname{int}(J)$ with first-$q$
separation above $4c_{\alpha,q}m^{-1/q}$ between different $U_i$ and
$\rho(f(u),f(\iota z_i))<\alpha/64$ on $U_i$, and pick
$p_i\in U_i\cap W$ by density of $W$.
\end{proof}

\begin{lemma}[Implantation; density of $\mathcal U_N(W_\bullet)$]\label{lem:density}
For every $N\ge1$, every $f\in\F(\Q)$, and every $\eta>0$ there exists
$g\in\mathcal U_N(W_\bullet)$ with $D(f,g)<\eta$.
\end{lemma}

\begin{proof}
Set $\alpha=\min(\eta/8,1/8)$ and $\theta=\alpha/8\le1/64$.  Let
$K_n=\max\bigl(N,q_n,\lceil\log_2(1/\alpha)\rceil+1\bigr)$.  All sites below are
$(\alpha,K_n)$-boxes \emph{anchored} in the sense that their restricted intervals lie
in $[2\theta,\,1-2\theta]$.

Let $J_1$ be any anchored $(\alpha,K_1)$-box.  Assume the anchored site $J_n$ is
defined.  Apply Lemma~\ref{lem:selection} to $(J_n,W_n,f_n,q_n,M_n)$ to obtain points
$p_{n,1},\dots,p_{n,M_n}\in\operatorname{int}(J_n)\cap W_n$ with first-$q_n$ pairwise
separation above $4c_{\alpha,q_n}M_n^{-1/q_n}$ and $f_n$-image diameter below
$\alpha/16$.  Choose pairwise disjoint open sets
$O_{n,i}\subseteq\operatorname{int}(J_n)\cap W_n$ with $p_{n,i}\in O_{n,i}$, first-$q_n$
separation above $c_{\alpha,q_n}M_n^{-1/q_n}$ between different $O_{n,i}$, and
$\rho(f_n(x),f_n(p_{n,i}))<\alpha/32$ for $x\in O_{n,i}$.  Since finite-support boxes
form a local base, choose bricks
$B_{n,i}\subseteq O_{n,i}$ with $p_{n,i}\in\operatorname{int}B_{n,i}$, support level
$R_n\ge K_n$, and $\diam_\rho(B_{n,i})<\gamma_n$.  As
$B_{n,i}\subseteq J_n$ in coordinates $\ell\le K_n$, the intervals $I^{n,i}_\ell$ for
$\ell\le N$ lie in the anchored site intervals, hence in $[2\theta,1-2\theta]$:
condition (S0) of Definition~\ref{def:scheme} holds, and (S1), (S2) and the window
membership hold by construction.

Let $c_{n+1}=f_n(p_{n,1})$ and define the next anchored site
$J_{n+1}=\prod_{\ell\le K_{n+1}}[u_\ell,u_\ell+\alpha]\times[0,1]^{\{K_{n+1}+1,\dots\}}$
with $u_\ell=\min\bigl(\max((c_{n+1})_\ell-\alpha/2,\,2\theta),\,1-2\theta-\alpha\bigr)$.
Every point $y\in J_{n+1}$ satisfies
$|y_\ell-(c_{n+1})_\ell|\le\alpha+2\theta$ for $\ell\le K_{n+1}$, hence
$\rho(y,c_{n+1})\le\alpha+2\theta+2^{-K_{n+1}}$.

Define the model map $q_{n,i}\in C(\Q,\Q)$ coordinatewise.  Writing
$I^{n,i}_\ell=[b_\ell,c_\ell]$ and
$\lambda_\ell(t)=\min\bigl(\max\bigl(\tfrac{t-b_\ell}{c_\ell-b_\ell},0\bigr),1\bigr)$:
\[
\pi_\ell q_{n,i}(x)=
\begin{cases}
(u_\ell-2\theta)+(\alpha+4\theta)\,\lambda_\ell(x_\ell), & \ell\le N,\\[2pt]
u_\ell+\alpha/2, & N<\ell\le K_{n+1},\\[2pt]
(c_{n+1})_\ell, & \ell>K_{n+1}.
\end{cases}
\]
The values in the first line range over $[u_\ell-2\theta,\,u_\ell+\alpha+2\theta]
\subseteq[0,1]$ by the anchoring of $J_{n+1}$.  Choose Urysohn functions
$\varphi_{n,i}:\Q\to[0,1]$ with $\varphi_{n,i}=1$ on $B_{n,i}$ and
$\varphi_{n,i}=0$ off $O_{n,i}$, and set
\[
g_n=\Bigl(1-\sum_i\varphi_{n,i}\Bigr)f_n+\sum_i\varphi_{n,i}\,q_{n,i},
\]
a coordinatewise convex combination, hence a continuous self-map of $\Q$ agreeing with
$f_n$ off $\bigcup_iO_{n,i}$ and with $q_{n,i}$ on $B_{n,i}$.

\emph{Error estimate.}  For $x\in O_{n,i}$ and $\ell\le N$, the value
$\pi_\ell q_{n,i}(x)$ lies in the $2\theta$-enlargement of the site interval
$[u_\ell,u_\ell+\alpha]$, every point of which is within $\alpha+4\theta$ of
$(c_{n+1})_\ell$; for $N<\ell\le K_{n+1}$ the value is the site centre, within
$\alpha/2+2\theta$ of $(c_{n+1})_\ell$; and for $\ell>K_{n+1}$ it equals
$(c_{n+1})_\ell$ exactly.  Hence
\[
\rho\bigl(q_{n,i}(x),c_{n+1}\bigr)
\le\sum_{\ell\le K_{n+1}}2^{-\ell}\bigl(\alpha+4\theta\bigr)
\le\alpha+4\theta .
\]
Also
$\rho(c_{n+1},f_n(x))\le\rho(f_n(p_{n,1}),f_n(p_{n,i}))+\rho(f_n(p_{n,i}),f_n(x))
<\alpha/16+\alpha/32<\alpha/8$.  Since $4\theta=\alpha/2$, we get
$\rho(g_n(x),f_n(x))\le\rho(q_{n,i}(x),f_n(x))<\alpha+\alpha/2+\alpha/8<2\alpha$,
and therefore $D(f,g)\le2\alpha<\eta$.

\emph{Verification of (S3).}  The bricks at time $n+1$, constructed at the next step
inside $\operatorname{int}(J_{n+1})$, have $\ell$-intervals contained in
$[u_\ell,u_\ell+\alpha]$ for every $\ell\le K_{n+1}$, so
$m^-_{n+1,\ell}\ge u_\ell$ and $m^+_{n+1,\ell}\le u_\ell+\alpha$ for $\ell\le N$.
If $x\in B_{n,i}$ has $x_\ell=b_\ell$, then
$\pi_\ell g_n(x)=\pi_\ell q_{n,i}(x)=u_\ell-2\theta\le m^-_{n+1,\ell}-2\theta$,
and on the opposite face the value is
$u_\ell+\alpha+2\theta\ge m^+_{n+1,\ell}+2\theta$.  Thus (S3) holds with margin
$2\theta>\theta$, for \emph{whatever} bricks the next step selects inside $J_{n+1}$.
Carrying out the induction for all $n$ yields $g\in\mathcal U_N(W_\bullet)$ with
$D(f,g)<\eta$.
\end{proof}

\section{The residual depth hierarchy}\label{sec:residual}

\begin{theorem}[Finite-depth genericity]\label{thm:main}
For every $N\ge1$ and every window sequence $W_\bullet$ the class
$\mathcal U_N(W_\bullet)$ is open and dense in $(\F(\Q),D)$.  Hence
\[
\mathcal S(W_\bullet)=\bigcap_{N\ge1}\mathcal U_N(W_\bullet)
\]
is residual.  Moreover, for any countable family $\{W^a_\bullet\}_{a\ge1}$ of window
sequences, $\bigcap_{a\ge1}\mathcal S(W^a_\bullet)$ is residual, and every
$g$ in this set admits, for every $a$ and every depth $N$, a strict depth-$N$
$W^a_\bullet$-windowed brick scheme with the exact one-step covering property of
Theorem~\ref{thm:covering}.
\end{theorem}

\begin{proof}
Openness is Lemma~\ref{lem:open}, density is Lemma~\ref{lem:density}, and
$(\F(\Q),D)$ is Baire by Lemma~\ref{lem:baire}; countable intersections of residual
sets are residual.  The covering property is Theorem~\ref{thm:covering}.
\end{proof}

\section{Generic entropy bursts in every tail}\label{sec:burst}

The depth hierarchy of Theorem~\ref{thm:main} is a statement about one-step structure.
We now prove that a genuinely \emph{dynamical}, infinite-horizon shadow of the dense
theorem of \cite{Shmalo2026} also holds residually.  Recall the tail Bowen data: for
$g\in\F(\Q)$ and $r\ge0$, $G^0_r=\operatorname{id}$,
$G^j_r=g_{r+j}\circ\cdots\circ g_{r+1}$, the Bowen metric
$\rho^{g}_{r,N}(x,y)=\max_{0\le j<N}\rho(G^j_r(x),G^j_r(y))$, and
$s_{r,N}(g,\varepsilon)$, the maximal cardinality of an
$(r,N,\varepsilon)$-separated set.  The tail scale entropy is
$H_{g,r}(\varepsilon)=\limsup_N\frac1N\log s_{r,N}(g,\varepsilon)$.

\begin{lemma}\label{lem:sep-open}
For fixed $r\ge0$, $N\ge1$, $\varepsilon>0$ and $m\in\mathbb N$, the set
$\{g\in\F(\Q):s_{r,N}(g,\varepsilon)\ge m\}$ is open.
\end{lemma}

\begin{proof}
Let $g$ admit an $(r,N,\varepsilon)$-separated set $E=\{x_1,\dots,x_m\}$, with
$\min_{u\ne v}\rho^g_{r,N}(x_u,x_v)=\varepsilon+2\kappa$ for some $\kappa>0$.  The maps
$g_{r+1},\dots,g_{r+N-1}$ are uniformly continuous; hence there is $\delta>0$ such
that $D(h,g)<\delta$ implies
$\sup_x\rho\bigl(H^j_r(x),G^j_r(x)\bigr)<\kappa$ for $0\le j<N$, by a finite induction
on $j$:
$\rho(H^{j}_r(x),G^{j}_r(x))\le
\rho\bigl(h_{r+j}(H^{j-1}_rx),g_{r+j}(H^{j-1}_rx)\bigr)
+\rho\bigl(g_{r+j}(H^{j-1}_rx),g_{r+j}(G^{j-1}_rx)\bigr)
\le D(h,g)+\omega_{g_{r+j}}\bigl(\rho(H^{j-1}_rx,G^{j-1}_rx)\bigr)$,
where $\omega_u$ is a modulus of continuity of $u$.  Then $E$ is
$(r,N,\varepsilon)$-separated for every such $h$.
\end{proof}

\begin{theorem}[Residual entropy bursts]\label{thm:burst}
The set
\[
\mathcal R \;=\;
\Bigl\{\,g\in\F(\Q):\ \forall r\ge0\ \forall q\ge1\ \forall j\ge1:\
\limsup_{N\to\infty}\ \sup_{0<\varepsilon<1/j}
\frac{\log s_{r,N}(g,\varepsilon)}{N\,\log(1/\varepsilon)}\;\ge\;q\Bigr\}
\]
is residual in $(\F(\Q),D)$.  In particular, for a residual set of systems,
\emph{every} tail exhibits separated orbit growth of every order $q$ at arbitrarily
fine scales:
$\limsup_{N}\sup_{\varepsilon<1/j}\frac{\log s_{r,N}}{N\log(1/\varepsilon)}=\infty$
for every $r$ and $j$.
\end{theorem}

\begin{proof}
For $r,q,j,N_0$ let
\[
V^{r,q,j}_{N_0}
=\bigcup_{\substack{\varepsilon\in\mathbb Q,\ 0<\varepsilon<1/j}}
\ \bigcup_{N\ge N_0}
\bigl\{\,g:\ s_{r,N}(g,\varepsilon)>\varepsilon^{-qN}\,\bigr\}.
\]
Each $V^{r,q,j}_{N_0}$ is open by Lemma~\ref{lem:sep-open}.  We claim it is dense.
Let $f\in\F(\Q)$ and $\eta>0$.  By the dense Hilbert-cube theorem of
\cite[Theorem~3.6]{Shmalo2026} there is $g$ with $D(f,g)<\eta$ and constants
$C_{g,q+1}$ such that
$H_{g,r}(\varepsilon)\ge(q+1)\log(1/\varepsilon)-C_{g,q+1}$
for every $r\ge0$ and all sufficiently small $\varepsilon>0$.  Choose a rational
$\varepsilon<1/j$ small enough that
$(q+1)\log(1/\varepsilon)-C_{g,q+1}>q\log(1/\varepsilon)$, i.e.\
$\log(1/\varepsilon)>C_{g,q+1}$.  Since
$H_{g,r}(\varepsilon)>q\log(1/\varepsilon)$, there are arbitrarily large $N$ with
$\frac1N\log s_{r,N}(g,\varepsilon)>q\log(1/\varepsilon)$, in particular some
$N\ge N_0$; then $s_{r,N}(g,\varepsilon)>\varepsilon^{-qN}$ and $g\in V^{r,q,j}_{N_0}$.
Thus each $V^{r,q,j}_{N_0}$ is open dense, and
$\mathcal R\supseteq\bigcap_{r,q,j,N_0}V^{r,q,j}_{N_0}$
is residual: indeed, if $g$ lies in the intersection then for all $r,q,j$ and every
$N_0$ there are $N\ge N_0$ and $\varepsilon<1/j$ with
$\log s_{r,N}(g,\varepsilon)>qN\log(1/\varepsilon)$, which is exactly the displayed
$\limsup\sup$ condition.
\end{proof}

\begin{remark}
Theorem~\ref{thm:burst} is the residual shadow of the dense
metric-mean-dimension theorem: it asserts unbounded \emph{burst} complexity, where
the certifying scale $\varepsilon$ may vary with the horizon $N$.  Upgrading it to
residual infinite \emph{upper metric mean dimension} --- where a single scale sequence
$\varepsilon\downarrow0$ must work for fixed-$\varepsilon$ horizon limits --- runs
precisely into the scale-pinning issue discussed in Section~\ref{sec:intro}; we leave
it as part of Problem~\ref{prob:main}.  Note also that near contracting systems the
fixed-scale sets $\{g:H_{g,r}(\varepsilon_0)>0\}$ fail to be dense (a $\delta$-ball
around a constant system has all tail orbits eventually $2\delta$-confined, killing
all separation at scales above $4\delta$), so the burst formulation, with its
scale flexibility, is the natural residual statement at this level of generality.
\end{remark}

\section{The obstruction: exact schemes have empty interior}\label{sec:obstruction}

We recall the exact class of \cite{Shmalo2026}.

\begin{definition}[{Exact windowed Hilbert horseshoe; \cite[Definition~12.1]{Shmalo2026}}]\label{def:exact}
A system $g\in\F(\Q)$ has an exact $W_\bullet$-windowed Hilbert horseshoe if there are
$\alpha\in(0,1/4)$, integers $K_n\ge q_n$, $(\alpha,K_n)$-boxes $J_n$, and
finite-support boxes $B_{n,1},\dots,B_{n,M_n}\subseteq\operatorname{int}(J_n)\cap W_n$
with $\diam_\rho(B_{n,i})<\gamma_n$, pairwise separation
$\dist_{q_n}(B_{n,i},B_{n,j})>c_{\alpha,q_n}M_n^{-1/q_n}$, and
\[
g_n(B_{n,i})=J_{n+1}\qquad(n\ge1,\ 1\le i\le M_n).
\]
Let $\EX$ denote the set of such systems.
\end{definition}

\begin{theorem}[Empty interior]\label{thm:empty}
$\EX$ is dense in $(\F(\Q),D)$ and has empty interior.  More generally, let
$\mathcal P\subseteq\F(\Q)$ be any class such that every $g\in\mathcal P$ admits some
$n\ge1$ and a compact set $A\subseteq\Q$ for which
$\pi_j\bigl(g_n(A)\bigr)=[0,1]$ for infinitely many coordinates $j$.  Then
$\mathcal P$ has empty interior.
\end{theorem}

\begin{proof}
Density of $\EX$ is \cite[Proposition~12.2]{Shmalo2026}.  For the rest, fix any
$g\in\F(\Q)$ and $\delta>0$; we produce $h$ with $D(h,g)<\delta$ and
$h\notin\mathcal P$.  Choose $L$ with $2^{-L}<\delta$ and define the
\emph{coordinate squeeze} $\Phi:\Q\to\Q$,
\[
\Phi(y)_j=
\begin{cases}
y_j, & j\le L,\\
\tfrac18+\tfrac34\,y_j, & j>L .
\end{cases}
\]
$\Phi$ is continuous and
$\rho(\Phi(y),y)=\sum_{j>L}2^{-j}\bigl|\tfrac18-\tfrac{y_j}{4}\bigr|
\le\tfrac38\,2^{-L}<\delta$.
Set $h=(\Phi\circ g_n)_{n\ge1}$, so $D(h,g)<\delta$.  For every $n$ and every
$x\in\Q$ and every $j>L$ we have $\pi_j h_n(x)\in[\tfrac18,\tfrac78]$.  Hence for
every $n$ and every compact $A$, $\pi_j(h_n(A))\subseteq[\tfrac18,\tfrac78]\ne[0,1]$
for all $j>L$: the image projections of $h$ are full in at most finitely many
coordinates.  Thus $h\notin\mathcal P$, and $\mathcal P$ contains no $D$-ball.

Finally, $\EX$ is such a class: if $g\in\EX$, then
$g_1(B_{1,1})=J_2$, and the $(\alpha,K_2)$-box $J_2$ has
$\pi_j(J_2)=[0,1]$ for every $j>K_2$, which is an infinite set of coordinates.
Hence $\operatorname{int}(\EX)=\emptyset$.
\end{proof}

\begin{remark}
Theorem~\ref{thm:empty} sharpens the observation of \cite[Section~14]{Shmalo2026}
from ``not asserted to be open'' to ``cannot meet any open set in a stable way'': the
dense class $\EX$ is topologically thin in the strongest possible sense compatible
with density.  The squeeze perturbation also destroys surjectivity of each map and
every property forcing deep full-range images, e.g.\ exactness onto any
finite-support box family with unbounded supports.  Together with
Remark~\ref{rem:depth-cost}, this closes the open-versus-dense question raised in
\cite[Section~14]{Shmalo2026}: openness is impossible at infinite depth, holds with
radius $\theta2^{-N-1}$ at depth $N$, and the residual object produced by Baire
fusion is the full depth hierarchy $\mathcal S(W_\bullet)$ of
Theorem~\ref{thm:main}.
\end{remark}

\section{What remains open}\label{sec:open}

\begin{problem}\label{prob:main}
Let $W_\bullet$ be a sequence of dense open subsets of $\Q$ and consider the class
$\mathcal C(W_\bullet)$ of systems $g\in\F(\Q)$ such that for every $r\ge0$ the tail
$\operatorname{tail}_rg$ admits a compact symbolic semifactor onto a full
variable-alphabet shift, a continuum-sized DC1-scrambled set visiting
$W_{r+j+1}$ at tail time $j$, infinite topological entropy, and
$H_{g,r}(\varepsilon)\ge q\log(1/\varepsilon)-C_{g,q}$ for every finite $q$.  By
\cite[Theorem~3.6]{Shmalo2026}, $\mathcal C(W_\bullet)$ is dense.  Is
$\mathcal C(W_\bullet)$ residual?  Is the subclass with infinite upper metric mean
dimension in every tail residual?
\end{problem}

The results of this paper bound the possible strategies.  By
Theorem~\ref{thm:empty}, no proof can proceed through an open class of systems with
unbounded-depth exact or margin-protected covering, since every such class has empty
interior; by Theorem~\ref{thm:main} the natural residual envelope is the depth
hierarchy $\mathcal S(W_\bullet)$, on which one-step covering is exact at every
finite depth but infinite-horizon orbit selection is not automatic; and by
Theorem~\ref{thm:burst} the burst form of infinite-order complexity is already
residual.  What is missing is exactly a mechanism producing, for a residual set of
systems, orbits threading infinitely many prescribed bricks --- equivalently, a
replacement for the backward induction of \cite[Lemma~13.1]{Shmalo2026} that consumes
only finite-depth information per step.  The quantitative erosion estimate of
Remark~\ref{rem:depth-cost} suggests that such a mechanism, if it exists, must
interleave depth and horizon nontrivially; we intend to return to this elsewhere.

\subsection*{Acknowledgments}
The research presented in this paper was supported by the European Research Council
(ERC) under the European Union's Horizon Europe research and innovation programme
(grant agreement No.\ 101041711), by the Israel Science Foundation (grant number
2258/19) and by the Israel Science Foundation (ISF Grant 4101/25).

\subsection*{Declaration on generative AI}
During the preparation of this work the author used generative AI tools to accelerate
drafting and revision.  All mathematical claims, proofs, and bibliographic information
were subsequently reviewed and edited by the author, who takes full responsibility for
the content of the manuscript.

\begin{thebibliography}{9}

\bibitem{Shmalo2026}
Y.~Shmalo,
\emph{Strict cubical horseshoes and all-tail full complexity in non-autonomous
dynamics}, manuscript, submitted, 2026.

\bibitem{Miranda}
C.~Miranda,
\emph{Un'osservazione su un teorema di Brouwer},
Boll.\ Un.\ Mat.\ Ital.\ (2) \textbf{3} (1940--41), 5--7.

\bibitem{Kulpa}
W.~Kulpa,
\emph{The Poincar\'e--Miranda theorem},
Amer.\ Math.\ Monthly \textbf{104} (1997), no.~6, 545--550.

\bibitem{KolyadaSnoha}
S.~Kolyada and L.~Snoha,
\emph{Topological entropy of nonautonomous dynamical systems},
Random Comput.\ Dynam.\ \textbf{4} (1996), 205--233.

\bibitem{BalibreaProblems}
F.~Balibrea,
\emph{On problems of topological dynamics in non-autonomous discrete systems},
Appl.\ Math.\ Nonlinear Sci.\ \textbf{1} (2016), no.~2, 391--404.

\bibitem{Oxtoby}
J.~C.~Oxtoby,
\emph{The Banach--Mazur game and Banach category theorem},
in: Contributions to the Theory of Games, Vol.~III,
Ann.\ of Math.\ Stud.\ \textbf{39}, Princeton Univ.\ Press, 1957, 159--163.

\bibitem{LindenstraussWeiss}
E.~Lindenstrauss and B.~Weiss,
\emph{Mean topological dimension},
Israel J.\ Math.\ \textbf{115} (2000), 1--24.

\bibitem{CRV}
M.~Carvalho, F.~B.~Rodrigues and P.~Varandas,
\emph{Generic homeomorphisms have full metric mean dimension},
Ergodic Theory Dynam.\ Systems \textbf{42} (2022), no.~1, 40--64.

\end{thebibliography}

\end{document}
